\section{The I and Subjectivity}

The common misconception is that the spirit exists in the body. However, the spirit is no other than Atma which is the principle of all states of the being at all degrees of its manifestation. Since all things are contained in their principle, the mind and body are contained in spirit, not the other way around.

Jivatma is Atma considered in relation to the human individuality; it resides at the center of this individuality and is designated symbolically by the heart. However, Atma can be neither manifested nor individualized nor embodied. As jivatma it appears to be individualized and embodied from the point of view of the individual human manifestation.

Analogously, for the ordinary man whose consciousness resides in the corporeal modality of the individual, his subtle modality, or center of subjective experience, likewise is believed to reside in the body. But the man who has realized his integral individuality undergoes a reversal and understands the body to be ``in'' the soul.

So when a being passes to supra-individual realization, he will then grasp the proper relationship between Atma and the individual states, that is, the latter appear as contained in the former.

In our ordinary state, we experience the I as jivatma at the centre of our subjective experience. This is characteristic of the individual state. However, this appearance is not the whole story. This realization is the awakened state of oneself as spirit, which is non-formal manifestation. As such, it is superior to the dream state which corresponds to subtle manifestation and the deep sleep state which corresponds to corporeal or gross manifestation. The I appears here to be individual.

However, beyond those three states of Atma, there is the fourth state, there is the unmanifested. This I then cannot be multiple, since quantity makes no sense in the unmanifested. The apparent reality of multiple I's or subjectivities at the third (or lower) state is an illusion.

Note especially that at this point the ego is transcended. So any accusations of egoism, the desire to control or dominate others is beside the point. Neither pride nor humility are applicable here, since both pertain to the sentimental and individual character. Similarly, this surpasses any notion of altruism, for where there is no I (in the relative sense), there is no longer an ``other''. As Guenon writes:

\begin{quotex}
This is a domain where all beings are one, `fused but not confused' according to Eckhart's expression and thereby they truly realize the words of Christ, `That they may be one even as the Father and I are one'.
\end{quotex}


To illustrate this point, Evola brings up the experience of dreaming, something that Guenon also does. In a dream, there appear to be many characters, each with its own centre of consciousness. However, on awakening we see that the entire dream, and all those characters, were simply the productions of our own mind. A similar realization occurs with the awakening to the fourth state.

References: (both by Rene Guenon)

\textit{Is the Spirit in the Body or the Body in the Spirit?}

\textit{Intellectual Pride}


\flrightit{Posted on 2011-02-26 by Cologero}

